\documentclass[11pt]{article}
\usepackage[T1]{fontenc}
\usepackage{lmodern,amsmath,amssymb,amsthm,mathtools}
\usepackage[margin=1in]{geometry}
\usepackage[hidelinks]{hyperref}
\newtheorem{theorem}{Theorem}[section]
\newtheorem{proposition}[theorem]{Proposition}
\newtheorem{lemma}[theorem]{Lemma}
\newcommand{\R}{\mathbb R}
\newcommand{\C}{\mathbb C}
\newcommand{\E}{\mathbb E}
\newcommand{\Pp}{\mathbb P}
\newcommand{\code}[1]{{\small\nolinkurl{#1}}}
\title{Brownian clock mixtures:\protect\\exact transforms, singular time evolution\protect\\and the limits of one-time information}
\author{Project Stardust\\\small A classical formalization companion}
\date{Reviewed editorial draft, October 30, 2026}
\begin{document}
\maketitle
\begin{abstract}
Two-sided Brownian motion evaluated at a signed Brownian clock and
Brownian motion evaluated at its reflected clock have the same one-time
Gaussian mixture laws. We reconstruct their exact characteristic
functions, the boundary term in their fourth-order time equation,
and the square-root cusp that prevents a finite initial time derivative
at any nonzero Fourier frequency. The laws are genuine probabilities,
but do not form a convolution semigroup. A separate covariance
calculation proves the reflected-clock observation is not a martingale
in its observed past. A hidden-clock local-time formula concerns the
scenery average and does not transfer to the unaveraged observation.
For arbitrary fixed even nonnegative compact smoothing profiles, an exact
Gaussian Stein decomposition recovers the original curvature boundary
correction in ordinary clock $L^2$, independently of the profile.
All proofs retain the standard Brownian variance convention and are
matched to a frozen Lean library. The Brownian-time PDE mechanism is
classical, with a direct primary comparison to Allouba and Zheng;
no new PDE, fractional generator, full observed-filtration calculus,
or priority claim is made.
\end{abstract}

\section{The two actual models and what is being compared}
Let $W$ be a standard Brownian motion on $[0,\infty)$ and let $B$
be an independent continuous two-sided Brownian scenery, normalized by
\begin{equation}\label{eq:scenery}
 \E B(r)=0,\qquad
 \E B(r)B(q)=\frac{|r|+|q|-|r-q|}{2}.
\end{equation}
In particular $B(r)$ has variance $|r|$. On each half-line the
scenery has the usual Brownian covariance, and the two halves have
independent Gaussian finite-dimensional laws. Define
\[
 Z_t^-=B(W_t),\qquad Z_t^+=B(|W_t|).
\]
The superscripts label signed and reflected clocks; they are not
positive and negative parts of a single random variable. Both
processes are constructed on the product of the independent source
spaces. Evaluation of continuous paths at a continuous random time
gives continuous sample paths. Joint measurability follows from the
jointly continuous evaluation map on continuous-path space and the
measurability of the two source paths. Thus the expectations below
refer to actual measurable compositions, not formal random times.

We prove results about their common one-time law $\mu_t$, and state
separately which assertions require the actual temporal coupling.
Equality of one-time laws supplies neither equality of process laws
nor equality of natural filtrations. The standard Brownian convention
here is $\E W_t^2=t$, with spatial generator $\tfrac12\partial_x^2$.
It differs from the variance-$2t$ convention used for the signed
higher-heat kernels in another companion.

The exact source construction is in \code{SignedIBMRealization}.
Its conditional Gaussian theorem retains the entire fixed clock path;
the one-time comparison is in \code{IBMOneTimeMixture}.
No change to either model or its verified hypotheses is made here.

\section{The common Gaussian mixture}
Write $\gamma_v=N(0,v)$, including $\gamma_0=\delta_0$.
The map $r\mapsto\gamma_{|r|}$ is a probability kernel.
For example, for bounded continuous $f$, its integral is
$\E f(\sqrt{|r|}G)$ with $G\sim N(0,1)$; dominated convergence
gives continuity in $r$. The usual generating-class argument then
gives measurability of the kernel on Borel sets. Define
\begin{equation}\label{eq:mixture}
 \mu_t(A)=\int_\R\gamma_{|r|}(A)\,\gamma_t(dr).
\end{equation}
This is a probability measure, since each conditional law has total
mass one and the mixing law is a probability.

\begin{theorem}\label{thm:law}
For every fixed $t\geq0$, both $Z_t^-$ and $Z_t^+$ have law $\mu_t$.
For every bounded Borel $f$,
\begin{equation}\label{eq:expectation}
 \E f(Z_t^-)=\E f(Z_t^+)
 =\int_\R\int_\R f(x)\,\gamma_{|r|}(dx)\,\gamma_t(dr).
\end{equation}
The same formula holds for nonnegative Borel tests, with possibly
infinite values, and for absolutely integrable tests.
\end{theorem}
\begin{proof}
Fix the clock source. At its value $W_t=r$, the signed observation
is the scenery coordinate $B(r)$, whose law is $N(0,|r|)$ by
\eqref{eq:scenery}. The reflected observation is $B(|r|)$, with
the same variance and law. Independence makes the scenery source
measure unchanged when the clock is fixed. Apply Tonelli first to
indicators of Borel sets. The resulting integral is exactly
\eqref{eq:mixture}, proving the two pushforward identities. Simple
approximation proves the nonnegative test formula. Applying it to
the absolute value, and then to positive and negative parts or real
and imaginary parts, proves every integrable version. In particular
bounded tests are integrable under these probability laws.
\end{proof}

There is a concrete limitation to this comparison. At fixed signed
outer times $1,-1$, the scenery covariance is zero, whereas at the
reflected times $|1|,|-1|$ it is one. More generally, conditional on
a clock path $w$, the signed two-time covariance is
\[
 \frac{|w(s)|+|w(t)|-|w(t)-w(s)|}{2},
\]
while the reflected covariance is $\min(|w(s)|,|w(t)|)$.
The formulas follow directly from \eqref{eq:scenery}; the latter
uses $a+b-|a-b|=2\min(a,b)$ for nonnegative $a,b$.
These are conditional scenery statements, not a claimed identification
of the unconditional two-time laws from one-dimensional marginals.

\subsection*{Moments and a one-time scaling check}
Let $G_1,G_2$ be independent standard normals on an auxiliary product
space. The mixture construction gives the distributional identity
\[
 Z_t^\pm\ \stackrel{\rm law}=\ t^{1/4}\sqrt{|G_2|}\,G_1.
\]
Indeed $W_t\stackrel{\rm law}=\sqrt t\,G_2$, and conditional
scenery variance is its absolute value. This identity is about each
fixed time. No simultaneous pathwise equality with this auxiliary
pair is asserted.
Gaussian absolute moments are finite of every fixed positive order,
so Tonelli proves integrability before any moment factorization.
Using $\E|G|=\sqrt{2/\pi}$, $\E G^2=1$,
$\E G^4=3$, and $\E G^8=105$ gives
\begin{equation}\label{eq:moments}
 \E Z_t^\pm=0,\qquad
 \E(Z_t^\pm)^2=\sqrt{\frac{2t}{\pi}},\qquad
 \E(Z_t^\pm)^4=3t,\qquad
 \E(Z_t^\pm)^8=315t^2.
\end{equation}
For example the eighth moment is $105\E W_t^4=315t^2$;
all conditional integrands are nonnegative. The normal moment values
follow by integration by parts against $e^{-x^2/2}$, with boundary
terms zero; the absolute first moment follows by integrating
$x e^{-x^2/2}$ on the positive half-line. For $t>0$ the kurtosis is
$3\pi/2>3$, so this nondegenerate one-time law is not Gaussian.
The displayed global moment evaluations and this diagnostic are human
consequences of the actual mixture; the formal actual quartic and
eighth-integrability constructions are retained in the source map.

\section{An exact transform and its boundary identity}
Put $c=(2\pi)^{-1/2}$ and
\[
 L(b)=\E e^{-b|G|}=2c\int_0^\infty e^{-r^2/2-br}\,dr,
 \qquad b\in\R.
\]
The integral is finite for every fixed $b$: completing the square
leaves a translated Gaussian tail. In the physical applications
below $b\geq0$, so $0<L(b)\leq1$ and $L(0)=1$.

\begin{proposition}\label{prop:transform}
For angular Fourier frequency $\xi\in\R$,
\begin{equation}\label{eq:characteristic}
 \widehat\mu_t(\xi):=\E e^{i\xi Z_t^\pm}
 =L\left(\frac{\xi^2\sqrt t}{2}\right).
\end{equation}
For $b\geq0$,
\begin{equation}\label{eq:boundary}
 \E\{|G|e^{-b|G|}\}=2c-bL(b),\qquad
 L'(b)=bL(b)-2c.
\end{equation}
The derivative at $b=0$ is an ordinary two-sided parameter derivative.
\end{proposition}
\begin{proof}
Conditional on the clock, the characteristic function of the scenery
normal is $e^{-\xi^2|W_t|/2}$. Its modulus is at most one, so Fubini
is justified directly on the source probability product. Brownian
scaling $|W_t|\stackrel{\rm law}=\sqrt t|G|$ gives
\eqref{eq:characteristic}.

For the boundary identity set $h_b(r)=e^{-r^2/2-br}$ on $r\geq0$.
It has derivative $h_b'=-(r+b)h_b$, value $h_b(0)=1$, and limit
zero at infinity. Both $h_b$ and $rh_b$ are integrable for $b\geq0$,
by domination by a Gaussian and its first moment. Integration of the
derivative on $[0,R]$ and then $R\to\infty$ yields
\[
 \int_0^\infty rh_b(r)\,dr
       =1-b\int_0^\infty h_b(r)\,dr.
\]
Multiply by $2c$. This is the first formula in \eqref{eq:boundary}.
Differentiating the probability integral for $L$ gives the negative
of this moment. On a neighborhood inside $b>0$, the derivative
integrands are dominated by $|G|$. Around zero, use instead
$|G|e^{|G|}$ for $|b|\leq1$. This has finite Gaussian integral:
its tail is a polynomial times $e^{-(|r|-1)^2/2}$ times a constant.
These majorants justify differentiation, including at zero.
\end{proof}

Completing the square also gives the useful exact expression
\[
 L(b)=2c e^{b^2/2}\int_b^\infty e^{-v^2/2}\,dv.
\]
This identity is a human re-expression of the defining integral;
it is not a separate mapped special-function implementation.
For $b>0$, the integral representation proves the elementary upper
bound $L(b)\leq2c/b$ by dropping $e^{-r^2/2}$ and integrating
$e^{-br}$. The bound $L(b)\leq1$ remains preferable near zero.

\section{The singular fourth-order time equation}
For $a\geq0$ define $H_a(t)=L(a\sqrt t)$ for $t\geq0$.
This is continuous at zero and is positive and at most one.
For $t>0$, the chain rule and \eqref{eq:boundary} give
\begin{equation}\label{eq:timeODE}
 H_a'(t)=\frac{a^2}{2}H_a(t)-\frac{ac}{\sqrt t}.
\end{equation}
At $a=0$ this is the constant identity. At $a>0$ the formula follows
from $\partial_t\sqrt t=1/(2\sqrt t)$ and
$aL'(a\sqrt t)=a^2\sqrt t L(a\sqrt t)-2ac$.

\begin{theorem}\label{thm:fourierODE}
At every $t>0$ and every real frequency,
\begin{equation}\label{eq:fourierODE}
 \partial_t\widehat\mu_t(\xi)
 =\frac{\xi^4}{8}\widehat\mu_t(\xi)
       -\frac{\xi^2}{2\sqrt{2\pi t}}.
\end{equation}
Equivalently, for $T\geq0$,
\begin{equation}\label{eq:integratedODE}
 \widehat\mu_T(\xi)
 =1+\frac{\xi^4}{8}\int_0^T\widehat\mu_s(\xi)\,ds
       -\xi^2 c\sqrt T.
\end{equation}
\end{theorem}
\begin{proof}
Substitute $a=\xi^2/2$ into \eqref{eq:timeODE} and use
\eqref{eq:characteristic}. To integrate to zero, first apply the
ordinary fundamental theorem on $[\varepsilon,T]$.
The bounded characteristic function is integrable in time and
$s^{-1/2}$ is integrable at zero. Letting $\varepsilon\downarrow0$
and using $\widehat\mu_0(\xi)=1$ proves
\eqref{eq:integratedODE}, since $\int_0^T s^{-1/2}ds=2\sqrt T$.
The $T=0$ identity is immediate.
\end{proof}

The sign of the fourth-order term is positive in Fourier space.
A stable biharmonic evolution instead has multiplier $e^{-t\xi^4}$.
Here the negative source in \eqref{eq:fourierODE}, depending on the
initial law, is indispensable. Omitting it does not preserve the actual
probability characteristic function. In fact the derivative is strictly
negative for $a>0$:
\[
 H_a'(t)=-\frac{a}{2\sqrt t}\E\{|G|e^{-a\sqrt t|G|}\}<0.
\]
Strictness follows because $|G|>0$ almost surely and the integrand is
integrable and positive there. Thus the positive fourth-order term
and the negative boundary term compensate in an exact probability
identity; the fourth-order term by itself is not its Markov generator.

\subsection*{A concrete spatial formulation with all interchanges checked}
Let $\phi\in L^1([-R,R];\C)$ for finite $R\geq0$, extended by zero, and set
\[
 f(x)=\int_{-R}^R\phi(\xi)e^{i\xi x}\,d\xi,
 \qquad u(t,x)=\E f(x+Z_t^\pm).
\]
Then $f$ is bounded and smooth, with every spatial derivative obtained
by inserting $(i\xi)^k$ under the integral. In particular
$\|f^{(k)}\|_\infty\leq R^k\|\phi\|_1$ for $k>0$.
Fubini uses the integrable majorant $|\phi(\xi)|$ and gives
\[
 u(t,x)=\int_{-R}^R\phi(\xi)e^{i\xi x}
                     \widehat\mu_t(\xi)\,d\xi.
\]
On a compact positive time interval $t\geq\delta>0$, the absolute
time derivative of the multiplier is at most
$R^4/8+R^2c/(2\sqrt\delta)$ by \eqref{eq:fourierODE}.
This constant times $|\phi|$ dominates time difference quotients
locally, and $R^4|\phi|$ dominates the fourth spatial derivative.
Consequently dominated differentiation proves
\begin{equation}\label{eq:PDE}
 \partial_tu(t,x)=\tfrac18\partial_x^4u(t,x)
                  +\frac{f''(x)}{\sqrt{8\pi t}},
 \qquad u(0,x)=f(x).
\end{equation}
Indeed $(i\xi)^4=\xi^4$ and $(i\xi)^2=-\xi^2$, giving both
the sign and the constant. Initial convergence is uniform in $x$:
dominate $|\phi(\xi)|\,|\widehat\mu_t(\xi)-1|$ by $2|\phi|$.
This compact-frequency spatial statement is an explicit human
Fubini consequence of the formal characteristic equation. It is not
an asserted general generator-domain theorem or a pathwise It\^o
formula. Its PDE mechanism is the classical Brownian-time equation
of Allouba--Zheng~\cite{az}, with the same Brownian normalization.

\section{The square-root cusp and failure of convolution composition}
The parameter derivative at zero yields
\begin{equation}\label{eq:cusp}
 \lim_{t\downarrow0}\frac{H_a(t)-1}{\sqrt t}=-2ca.
\end{equation}
This follows by substituting $b=a\sqrt t$ into the defining difference
quotient of $L'(0)$; $a=0$ is constant. For $a>0$, the limit is
strictly negative, so $H_a$ cannot have a finite right time derivative
at zero. If it had derivative $d$, then
$(H_a(t)-1)/\sqrt t=((H_a(t)-1)/t)\sqrt t\to0$,
contradicting \eqref{eq:cusp}. The formal source also proves this by
composing a hypothetical time derivative with $t=s^2$: the chain
rule gives derivative zero, whereas $L(as)$ has derivative $-2ca$.
Thus the cusp conclusion does not rest on a formal interchange of a
singular derivative at the origin.

\begin{proposition}\label{prop:nosemigroup}
The probability family $\mu_t$ is not a convolution semigroup.
In particular it cannot be the marginal family at zero of a process
with stationary independent increments.
\end{proposition}
\begin{proof}
Suppose $\mu_{s+t}=\mu_s*\mu_t$ for all nonnegative $s,t$.
Convolution multiplies characteristic functions, so at $\xi=1$
and $t=1$ we would have
\[
 H_{1/2}(s+1)=H_{1/2}(s)H_{1/2}(1),\qquad s\geq0.
\]
The last factor is strictly positive. Therefore
$H_{1/2}(s)=H_{1/2}(s+1)/H_{1/2}(1)$ would have a finite
right derivative at zero, since the numerator is differentiable at
positive time one. This contradicts the cusp.
For stationary independent increments starting at zero, the law at
$s+t$ is the convolution of the law at $s$ with that of the
independent length-$t$ increment, giving precisely the excluded law.
\end{proof}

There is a second explicit check using \eqref{eq:moments}.
The centered law $\mu_1*\mu_1$ has second moment
$2\sqrt{2/\pi}$, whereas $\mu_2$ has second moment $2/\sqrt\pi$.
Both moments are finite, and independent centered summands add
their variances, so the measures differ. This particular witness is
a human deduction from the actual mixture moments; the formal
convolution obstruction uses the preceding transform argument.
Neither argument excludes a general nontranslation Markov model
merely from its one-time marginals. Convolution composition and
the full Markov property have different hypotheses.

\section{A temporal obstruction requiring the actual reflected model}
We now use $Z_t=Z_t^+$ on its original source product, rather than
only the family $\mu_t$. All pairings are integrable by
\eqref{eq:moments} and Cauchy--Schwarz. Conditional on the clock,
\eqref{eq:scenery} at nonnegative times gives
\[
 \E_B\{Z_s(Z_t-Z_s)\}
  =\min(|W_s|,|W_t|)-|W_s|
  =-(|W_s|-|W_t|)_+.
\]
Fubini is justified by the original joint $L^2$ bounds. Hence
\begin{equation}\label{eq:negative}
 \E\{Z_s(Z_t-Z_s)\}
      =-\E (|W_s|-|W_t|)_+.
\end{equation}
The nonnegative right-hand integrand is bounded by $|W_s|$, so is
integrable. At $s=1,t=2$ its expectation is strictly positive.
Indeed the event
\[
 1<W_1<2,\qquad -1<W_2-W_1<-\tfrac12
\]
has positive probability: its two coordinates are independent
nondegenerate standard normals, each interval has positive Gaussian
mass, and the event probability is their positive product. On it
$0<W_2<W_1$, so $(|W_1|-|W_2|)_+>0$.
Thus \eqref{eq:negative} is strictly negative at those times.

If $Z$ were a square-integrable martingale for its raw observed past,
then $\E(Z_2-Z_1\mid\sigma(Z_r:r\leq1))=0$.
Pair with the past-measurable $Z_1\in L^2$, or first with bounded
truncations and pass by Cauchy--Schwarz, to get zero in
\eqref{eq:negative}, a contradiction. More generally, every closed
past subspace containing $Z_1$ fails to project $Z_2$ onto $Z_1$:
the orthogonal projection identity paired with $Z_1$ would give
that same zero. This is the exact actual temporal obstruction;
it does not claim that the process has no other observed martingales,
or that a Markov property is disproved by the martingale test alone.

\section{A source mean has a hidden-clock local-time formula}
The conditional scenery average of a Fourier cosine is
\[
 M_t(\xi)=\E_B\cos(\xi B(|W_t|))
         =e^{-a|W_t|},\qquad a=\xi^2/2.
\]
This is an actual scalar function of the Brownian clock. Its ordinary
Brownian filtration and the observed IBM filtration are not identified
by the one-time mixture theorem.

Use the symmetric local time convention fixed by Tanaka's formula
\[
 |W_t|=\int_0^t\operatorname{sgn}(W_s)\,dW_s+L_t^0,
 \qquad \operatorname{sgn}(0)=0.
\]
The original constructed Brownian mass path is this $L^0$; its
continuous Stieltjes measure is supported where $W_s=0$.
Tanaka and that supported-mass construction are established in the
sticky-clock companion~\cite{sticky}; here we use their stated
classical Brownian identity.

For $a\geq0$ a genuine global $C^2$ extension of $x\mapsto e^{-ax}$
from the nonnegative half-line is
\[
 F_a(x)=\begin{cases} e^{ax}-2ax,&x<0,\\e^{-ax},&x\geq0.\end{cases}
\]
Its one-sided values and first derivatives agree at zero, with
$F_a(0)=1$ and $F_a'(0)=-a$. Its second derivative is
$a^2 e^{-a|x|}$ on both sides and continuously at zero.
The derivative is bounded by $2a$ and the second derivative by
$a^2$. Thus the ordinary $C^2$ semimartingale It\^o theorem applies
to $F_a(|W|)$ on each finite horizon. The quadratic variation of
$\int\operatorname{sgn}(W_s)dW_s$ is $t$: the zero set has
Lebesgue measure zero almost surely, by
$\E\int_0^T1_{\{W_s=0\}}ds=\int_0^T\Pp(W_s=0)ds=0$.
The finite variation local time contributes no quadratic variation.
The supported Stieltjes integral uses $F_a'(|W_s|)=-a$ when
$dL_s^0\ne0$. It\^o therefore gives, at every fixed $T$ almost surely,
\begin{equation}\label{eq:clockIto}
 M_T(\xi)-1
 =-a\int_0^T\operatorname{sgn}(W_s)e^{-a|W_s|}\,dW_s
       -aL_T^0+\frac{a^2}{2}\int_0^T M_s(\xi)\,ds.
\end{equation}
The stochastic integrand is bounded by $a$ and predictable in the
clock filtration, so its integral is square integrable with mean zero.
The time integral is bounded by $T$. Tanaka implies
$L_T^0=|W_T|-\int_0^T\operatorname{sgn}(W_s)dW_s\in L^2$,
and $\E L_T^0=\E|W_T|=2c\sqrt T$.
Taking expectations in \eqref{eq:clockIto} is consequently justified
and gives exactly \eqref{eq:integratedODE}.

The boundary term of the one-time PDE has a concrete hidden-clock
origin. This does not make \eqref{eq:clockIto} an It\^o formula for
the unaveraged $B(|W_t|)$ in its observed filtration. Its left side
has already integrated out the scenery, and its stochastic integral
is with respect to the original clock Brownian motion. The source
theorem is almost sure for each fixed horizon; no simultaneous
all-horizons version is claimed here.

\section{A profile-independent boundary correction in ordinary clock $L^2$}
The boundary in the preceding formula can also be recovered from the
literal curvature of spatially smoothed scenery. We identify the
residual in the original smoothing and prove convergence in ordinary
clock $L^2$. Averaging over scenery remains essential throughout
this section.

Fix a real Schwartz function $g$, a shift $a\in\R$, and an even,
nonnegative $C^2$ function $\rho$ of compact support with
$\int_\R\rho=1$. Choose $R>0$ with support in $[-R,R]$. For every
$\varepsilon>0$ set
\[
 \rho_\varepsilon(u)=\varepsilon^{-1}\rho(u/\varepsilon),\qquad
 S_\varepsilon(r)=\int_\R\rho_\varepsilon(u)B(r-u)\,du,
 \qquad v_\varepsilon(r)=\E_B S_\varepsilon(r)^2.
\]
The integral defining $S_\varepsilon$ is a Bochner integral in the
original scenery $L^2$. Compact support and continuity of
$r\mapsto B(r)$ in that Hilbert norm prove its integrability.
The result is a centered Gaussian variable: Riemann sums are
Gaussian and converge in $L^2$, hence their characteristic
functions converge to that of the centered Gaussian with the
limiting variance. Derivatives in $r$ are obtained by moving them
to the compact $C^2$ kernel. Write $S_\varepsilon'$ and
$S_\varepsilon''$ for these genuine Gaussian $L^2$ derivatives.
At each fixed $r$, their joint Gaussian moments of every finite
order exist.

\subsection*{The variance curvature is the original kernel}
The covariance \eqref{eq:scenery} and two applications of Fubini give
\begin{equation}\label{eq:smoothedvariance}
 v_\varepsilon(r)=\int\rho_\varepsilon(u)|r-u|\,du
 -\tfrac12\iint\rho_\varepsilon(u)\rho_\varepsilon(w)|u-w|\,du\,dw.
\end{equation}
All variables $u,w$ range over compact supports; the covariance
is bounded there for a fixed $r$, so the Fubini steps are absolute.
The second term is independent of $r$. Splitting the first
integral at $u=r$ and applying the fundamental theorem of calculus
gives
\begin{equation}\label{eq:variancecurvature}
 v_\varepsilon'(r)=2\int_{-\infty}^r\rho_\varepsilon(u)\,du-1,
 \qquad v_\varepsilon''(r)=2\rho_\varepsilon(r).
\end{equation}
Thus $|v_\varepsilon'|\leq1$, $v_\varepsilon'$ equals the sign of
$r$ when $|r|>\varepsilon R$, and $v_\varepsilon$ is even.
The curvature is twice the original profile, rather than its
square or self-convolution.

Jensen's inequality in the scenery Hilbert space gives the
uniform original-coupling estimate
\begin{equation}\label{eq:sceneryapprox}
 \E_B|S_\varepsilon(r)-B(r)|^2
 \leq\int\rho_\varepsilon(u)\E_B|B(r-u)-B(r)|^2du
 =\int\rho_\varepsilon(u)|u|du\leq\varepsilon R.
\end{equation}
In particular $v_\varepsilon(0)\leq\varepsilon R$. The slope
bound yields $v_\varepsilon(r)\leq2\varepsilon R$ when
$|r|\leq\varepsilon R$. Also \eqref{eq:smoothedvariance} gives
$|v_\varepsilon(r)-|r||\leq2\varepsilon R$ for every $r$:
the first-term error is at most $\varepsilon R$ and the absolute
second term at most $\varepsilon R$.
\par\noindent\code{IBMCompactVarianceGeometry.variance_second}.
\par\noindent\code{IBMCompactVarianceGeometry.slope_bound}.
\par\noindent\code{IBMDilatedMollifier.scenery_error}.

\subsection*{The literal Gaussian correction and its exact decomposition}
Define the scalar source-curvature residual
\begin{align}\label{eq:originalresidual}
 C_\varepsilon(r)
 &=\E_B\left[g''(a+S_\varepsilon(r))(S_\varepsilon'(r))^2
            +g'(a+S_\varepsilon(r))S_\varepsilon''(r)\right]
       -\tfrac14\E_B g^{(4)}(a+S_\varepsilon(r)).
\end{align}
These are actual integrals against the original scenery law.
Bounded Schwartz derivatives and finite Gaussian second moments
give absolute integrability of both terms inside the expectation.
No conditional projection of a singular limiting field defines them.

For any centered jointly Gaussian triple $(X,Y,V)$ and a smooth
function with the bounded derivatives needed below, Gaussian
integration by parts gives
\begin{align*}
 \E[Y^2h(X)]&=(\E Y^2)\E h(X)+(\E XY)^2\E h''(X),\\
 \E[Vh(X)]&=(\E XV)\E h'(X).
\end{align*}
For completeness, realize the triple as linear forms in a finite
standard Gaussian vector $\xi$. The elementary identity
$\E[\xi_jF(\xi)]=\E[\partial_jF(\xi)]$ gives the second formula
once. Apply it first to $Yh(X)$ and then to $h'(X)$ for the first
formula. Gaussian moments dominate every term. The argument
uses no inverse covariance matrix and includes degenerate triples.

Apply these identities with $X=S_\varepsilon$, $Y=S_\varepsilon'$
and $V=S_\varepsilon''$. Hilbert-space differentiation gives
\[
 v_\varepsilon'=2\E_B(S_\varepsilon S_\varepsilon'),\qquad
 v_\varepsilon''=2\E_B\bigl[(S_\varepsilon')^2+
                                      S_\varepsilon S_\varepsilon''\bigr].
\]
Using \eqref{eq:variancecurvature}, the literal residual is exactly
\begin{equation}\label{eq:residualdecomp}
 C_\varepsilon(r)=\rho_\varepsilon(r)A_{2,\varepsilon}(r)
       +\frac{(v_\varepsilon'(r))^2-1}{4}A_{4,\varepsilon}(r),
 \qquad A_{j,\varepsilon}(r)=\E_B g^{(j)}(a+S_\varepsilon(r)).
\end{equation}
The constants $1/4$ and $2$ come from the displayed actual
covariance derivatives. They are not adjusted to fit a desired limit.
\par\noindent\code{IBMCorrectionStein.correction_identity}.
\par\noindent\code{IBMCompactOriginalBoundary.originalResidual_literal}.
\par\noindent\code{IBMCompactOriginalBoundary.originalResidual_identity}.

\subsection*{Constructing the occupation mass and its ordinary $L^2$ limit}
Let $L_T^0$ denote the symmetric Tanaka local time of the original
clock, with the normalization already used above,
\[
 L_T^0=|W_T|-\int_0^T\operatorname{sgn}(W_s)\,dW_s.
\]
Evenness of $\rho$ makes the literal reflected occupation mass
\[
 M_\varepsilon(T)=\int_0^T\rho_\varepsilon(|W_s|)ds
                =\int_0^T\rho_\varepsilon(W_s)ds.
\]
Ordinary $C^2$ It\^o calculus applied to $v_\varepsilon(W)$ and
\eqref{eq:variancecurvature} gives, for each fixed $T$,
\begin{equation}\label{eq:massito}
 M_\varepsilon(T)=v_\varepsilon(W_T)-v_\varepsilon(0)
                   -\int_0^Tv_\varepsilon'(W_s)\,dW_s.
\end{equation}
The derivatives are bounded at each fixed width, and
$|v_\varepsilon(W_T)-v_\varepsilon(0)|\leq|W_T|$.
Thus every term is in clock $L^2$, and the triangle inequality
and It\^o isometry give the useful human bound
\begin{equation}\label{eq:massbound}
 \|M_\varepsilon(T)\|_2\leq2\sqrt T.
\end{equation}

The endpoint difference in \eqref{eq:massito} tends to $|W_T|$
in $L^2$, by the uniform variance estimate above. For every
$s>0$, $\Pp(W_s=0)=0$ by its Gaussian density; Tonelli therefore
gives zero $ds\,d\Pp$ measure to $W_s=0$ on $[0,T]$.
At every nonzero value, $v_\varepsilon'$ tends to the sign,
and both are bounded by one. Dominated convergence in
$L^2(ds\,d\Pp)$ and It\^o isometry prove
\begin{equation}\label{eq:masslimit}
 M_\varepsilon(T)\longrightarrow L_T^0
                    \quad\hbox{in clock }L^2\quad(\varepsilon\downarrow0).
\end{equation}
This constructs the needed mass limit without assuming a spatial
local-time density or a pathwise occupation formula. It holds
through all positive widths, not only a selected sequence.
One may verify this last assertion along every positive sequence
converging to zero, which characterizes this limit in a metric space.
\par\noindent\code{IBMCompactClockMass.clockMass_identity}.
\par\noindent\code{IBMCompactClockMass.clockMass_reflected_ae}.
\par\noindent\code{IBMDilatedClockMassLimit.reflected_pathIntegral_continuum_tendsto}.

\begin{theorem}[Original boundary correction]\label{thm:boundaryprofile}
For every fixed $T\geq0$ and every fixed profile satisfying the
stated hypotheses, the literal original-law random functional
\[
 H_\varepsilon(T)=\int_0^TC_\varepsilon(|W_s|)\,ds
\]
belongs to clock $L^2$ and satisfies
\begin{equation}\label{eq:boundarylimit}
 H_\varepsilon(T)\longrightarrow g''(a)L_T^0
        \quad\hbox{in ordinary clock }L^2
                         \quad(\varepsilon\downarrow0).
\end{equation}
The limit is independent of the fixed admissible profile.
\end{theorem}
\begin{proof}
Put $K=\|g'''\|_\infty$ and $N=\|g^{(4)}\|_\infty$.
On the support of $\rho_\varepsilon(r)$, the variance bound and
Cauchy--Schwarz give
\[
 |A_{2,\varepsilon}(r)-g''(a)|
       \leq K\E_B|S_\varepsilon(r)|\leq K\sqrt{2\varepsilon R}.
\]
Since $\rho_\varepsilon\geq0$, the first term of
\eqref{eq:residualdecomp}, integrated along the reflected clock,
differs from $g''(a)M_\varepsilon(T)$ by at most
$K\sqrt{2\varepsilon R}\,M_\varepsilon(T)$ pointwise.
Its $L^2$ error is at most
$K\sqrt{2\varepsilon R}\,\|M_\varepsilon(T)\|_2$, which tends
to zero by \eqref{eq:massbound}. Each such weighted integral
is in $L^2$ by this estimate and the mass identity.

The second residual term is bounded in absolute value by $N/4$.
It tends to zero for every nonzero $r$; in fact it vanishes
when $|r|>\varepsilon R$. It is a continuous function of $r$,
as follows from continuity of the Gaussian $L^2$ scenery curve
and the Lipschitz property of $g^{(4)}$. Composition with the
continuous clock is jointly measurable. The zero-occupation
fact above and dominated convergence on the finite product
measure $ds\,d\Pp$ show its time integral tends to zero in
clock $L^2$: Cauchy--Schwarz bounds the squared integral by
$T$ times the integral of its square. It is also in $L^2$
at each width, since its absolute time integral is at most $NT/4$.
Combine the two terms and \eqref{eq:masslimit} to obtain
\eqref{eq:boundarylimit}. The case $T=0$ is immediate.
\end{proof}
\par\noindent\code{IBMCompactSourceAverage.sourceAverage_layer_error}.
\par\noindent\code{IBMCompactWeightedBoundary.weightedClass_continuum_tendsto}.
\par\noindent\code{IBMCompactBoundaryRemainder.integralClass_continuum_tendsto}.
\par\noindent\code{IBMCompactOriginalBoundary.originalPathIntegral_continuum_tendsto}.

The original residual, the whole positive-width limit and its ordinary
clock $L^2$ topology are exactly the mapped theorem's scope. The
explicit uniform estimate \eqref{eq:massbound} is a human
consequence of the displayed actual It\^o identity and slope bound.
The profile is fixed while its width shrinks: no uniformity over
varying profiles with uncontrolled support radii is asserted.
Evenness is used to replace $W$ by $|W|$ in the occupation mass;
positivity supplies the weighted error and slope bounds.

There is no extra factor $1/2$ in \eqref{eq:masslimit}: the
reflected integral equals the full even-kernel integral. If the
curvature appears with the usual It\^o prefactor $1/2$, the
residual contribution is consequently $g''(a)L_T^0/2$.
The subtraction in \eqref{eq:originalresidual} concerns a
\emph{scenery expectation}. This theorem does not prove a joint
scenery/clock $L^2$ limit of the unaveraged curvature, convergence
of an intrinsic predictable IBM integrand, or a physical-time
It\^o formula for $g(a+Z_T^+)$. It recovers a precise original
boundary correction with all of those other obligations kept distinct.

\section{Primary comparison and exact formal scope}
Allouba and Zheng~\cite{az}, Theorem 0.1, equations (0.1)--(0.2),
already give the Brownian-time initial-data source and fourth-order
PDE. Their Section 1, first part of the proof, obtains the source
by integrating the reflected clock density by parts twice. We read
that theorem statement and the complete Brownian-time part of its
proof on printed pages 2 and 4--5 of the primary preprint. We do
not import the excursion-limit appendix, exit-time theorems or their
separate half-derivative conditional generator. In particular the
phrase ``half derivative'' there is not a license to assert a
Caputo equation or an observed-past representation theorem here.

The contributions of this companion are a readable reconstruction
and an exact correspondence to verified source statements. The
Gaussian mixture, Brownian-time PDE, normal moment tests, cusp and
convolution test are classical mechanisms; no priority claim follows
from their formalization or from grouping them in one paper.
The compact-frequency PDE, integrated endpoint formula, explicit
moment evaluations and particular second-moment convolution witness
are labelled human consequences with proofs, not newly named Lean
declarations. The primary literature already supplies a broader
PDE setting than the compact-frequency illustration.

The frozen baseline is
\begin{center}\small\ttfamily
bab688bb6140cd6316c80680095a914ea0d346db.
\end{center}
Principal source families, read in full, are
\code{IBMOneTimeMixture} and \code{IBMOneTimeFourier}.
\par\noindent
The boundary and time calculations use
\code{IBMGaussianLaplaceBoundary}.
\par\noindent
Their derivative chain is \code{IBMClockTransformDerivative} and
\code{IBMClockTimeEvolution}.
\par\noindent
The exact final characteristic equation is in
\code{IBMFourierTimeEquation}.
\par\noindent
The initial and convolution obstructions are in
\code{IBMInitialCusp} and \code{IBMConvolutionObstruction}.
\par\noindent
The actual backtracking test is in \code{IBMClockBacktracking}.
\par\noindent
The past-projection consequence is in
\code{IBMTemporalCovarianceObstruction}.
\par\noindent
The source-mean and damped reflection identities are in
\code{IBMOneTimeSourceMeanIto}.
\par\noindent
The original $C^2$ extension is in \code{GaussianDampedReflectedIto}.

The accompanying map gives exact declaration headers, source line
links and byte hashes; the baseline receipt checks the existing full
build and transitive axiom audit. The editorial work changes no
mathematical source and runs no redundant Lean build.
Nothing here completes intrinsic observed-filtration IBM calculus,
constructs a Markov semigroup from $\mu_t$, establishes a full
physical-time generator domain, or upgrades a one-time comparison
to an equality of process laws. Those distinctions are mathematical
hypotheses, not matters of notation.

\begin{thebibliography}{9}
\bibitem{az} H. Allouba and W. Zheng,
Brownian-time processes: the PDE connection and the half-derivative
generator, \emph{Ann. Probab.} \textbf{29} (2001), 1780--1795.
Primary preprint arXiv:1005.3801v1, deposited 2010;
manuscript dated January 2001.
\url{https://arxiv.org/abs/1005.3801}.
\bibitem{sticky} Project Stardust,
\emph{Brownian sticky clocks: construction, holding time and causal
recovery in the observed filtration}, reviewed editorial companion,
\code{manuscripts/sticky-clock-recovery/}. The Brownian local-time
and supported-mass construction is the prerequisite used in Section 7.
\end{thebibliography}
\end{document}
